%% Appendix A: the open exam question pool, transcribed verbatim, plus the
%% mapping from pool questions to chapters.
%% Source: material/pool/open-pool-questions.pdf (draft version).
%% When the pool is updated, update this file AND the mapping table below.

\chapter{The Open Exam Question Pool}
\label{app:pool}

The course is examined with the
\href{https://openpoolexams.science.uu.nl/}{open pool exam method}. At least
$30\%$ of the questions of the midterm and of the final exam are drawn
\emph{verbatim} from the pool reproduced below. The remaining questions test the
same material on instances you have not seen.

The pool is a draft and may change; the version transcribed here is the one in
\texttt{material/pool/open-pool-questions.pdf}. Where the pool numbers its
sources by lecture, we give the corresponding chapter of this book.

\section{Where each question is answered}
\label{sec:pool-mapping}

\begin{center}
\begin{tabular}{@{}llll@{}}
\toprule
\textbf{Q} & \textbf{Title} & \textbf{Pool lecture} & \textbf{Chapter} \\
\midrule
1  & Stable Matchings                    & 1  & \cref{ch:stable-marriage} \\
2  & Minimum Vertex Cut                  & 2  & \cref{ch:max-flow} \\
3  & Escaping a Grid                     & 2  & \cref{ch:max-flow} \\
4  & Phasing out capital equipment       & 3  & \cref{ch:min-cost-flow} \\
5  & Planar Graphs                       & 4  & \cref{ch:planar} \\
6  & Real Models of Computation          & 5  & \cref{ch:real-models} \\
7  & Partition into Cycles               & 6  & \cref{ch:matching} \\
8  & Existential Theory of the Reals     & 7  & \cref{ch:etr} \\
9  & Exact Algorithm for Exact 3-SAT     & 8  & \cref{ch:exact-i} \\
10 & Bisection Width                     & 8  & \cref{ch:exact-i} \\
11 & Duality and geometric $\ER$-hardness & 9  & \cref{ch:geometric-er} \\
12 & Inclusion/exclusion for set cover   & 10 & \cref{ch:exact-ii} \\
13 & Approximation non-symmetric TSP     & 10 & \cref{ch:approx-i} \\
14 & Approximating Steiner Tree          & 11 & \cref{ch:approx-ii} \\
15 & Parameterized complexity            & 12 & \cref{ch:fpt-i} \\
16 & Treewidth of grid graphs            & 13 & \cref{ch:treewidth} \\
17 & Treewidth and cliques               & 13 & \cref{ch:treewidth} \\
18 & L-reduction for Max Cut on Multigraphs & 14 & \cref{ch:hardness} \\
19 & Vertex Cover Kernel                 & 15 & \cref{ch:kernelization}, \cref{ch:randomized} \\
\bottomrule
\end{tabular}
\end{center}

\begin{remark}
Question~11 is titled \emph{Exact Algorithm for Exact 3-SAT} in the pool PDF,
which is a copy-paste error: its items are about point--line duality and
geometric $\ER$-completeness. We keep the numbering but use the accurate title.
\end{remark}

\section{The questions}
\label{sec:pool-questions}

\subsection*{Lecture 1: Stable Marriage}

\paragraph{Question 1. Stable Matchings.}
This set of questions captures the learning goals of the first lecture on the
stable marriage problem.
\begin{enumerate}
  \item Define an Euler Tour and show that every connected graph with all even
        degrees must have an Euler Tour.
  \item Define the term stable marriage problem, blocking pair and stable
        matching.
  \item Define the Greedy Algorithm as in the lecture and show that it might run
        forever.
  \item Explain the Gale Shapley algorithm and show that it returns a stable
        matching.
  \item Define a men optimal stable matching and show that the Gale Shapley
        algorithm returns a men optimal stable matching.
\end{enumerate}

\subsection*{Lecture 2: Min Flow Max Cut}

\paragraph{Question 2. Minimum Vertex Cut.}
Show that determining the maximum flow in a network with arc and vertex capacity
constraints can be reduced to an ordinary maximum-flow problem on a flow network
of comparable size.

\paragraph{Question 3. Escaping a Grid.}
An $n \times n$ grid is an undirected graph consisting of $n$ rows and $n$
columns of vertices. We denote the vertex in the $i$th row and $j$th column by
$v_{i,j}$. All vertices in a grid have exactly four neighbours, except for the
\emph{boundary} points, which are the vertices $v_{i,j}$ with $i = 1$, $i = n$,
$j = 1$ or $j = n$.

Given $m \le n^2$ starting points $v_{x_1,y_1}, \ldots, v_{x_m,y_m}$ in the
grid, the \emph{grid escape problem} is to determine if there are $m$
vertex disjoint paths from the starting points to any $m$ different points in
the boundary.

Describe an algorithm with running time polynomial in $n$ that solves the escape
problem and analyse its running time. (Be careful that $n$ and $m$ describe the
size of the grid and number of starting points, not the number of vertices and
arcs as usual.)

\subsection*{Lecture 3: Min Flow Max Cut}

\paragraph{Question 4. Phasing out capital equipment.}
A shipping company wants to phase out a fleet of (homogeneous) general cargo
ships over a period of $p$ years. Its objective is to maximize its cash assets at
the end of the $p$ years by considering the possibility to prematurely selling
ships and temporary replacing them by ships they rent.

The company faces a non-increasing demand for ships: every year, the number of
ships needed is at most the number needed the previous year. Let $d(i)$ be the
number of ships needed in year $i$. The company can possibly have more than
$d(i)$ ships in year $i$, but may not have less than these.

If a ship is sold in year $i$, this yields $s_i$ euros. Each ship the company has
in year $i$ gives $r_i$ euros as profit. If the company has less than $d(i)$
ships in year $i$, it must hire additional ships. This costs $h_i$ in year $i$.
Assume that $h_i > r_i$ for each $i$.

The company wants to meet its commitments (have sufficient ships either owned or
rented each year), and have as much cash assets as possible at the end of the
$p$th year.

Formulate this problem as a minimum cost flow problem.
(Exercise 9.6 from \emph{Network Flows}, Ahuja, Magnanti, Orlin.)

\subsection*{Lecture 4: Planar Graphs}

\paragraph{Question 5. Planar Graphs.}
This set of questions captures the learning goals of the lecture on the planar
graphs. We assume that $G$ is a planar graph on $n$ vertices and $m$ edges.
\begin{enumerate}
  \item Give and prove the Euler Formula.
  \item Show that $m \le 3n-6$ and the average degree below $6$.
  \item Define graph minors and state the Kuratowski theorem.
  \item Show that planar graphs are $5$-colorable.
  \item Show that $K_{3,3}$ is non-planar in an elementary way.
  \item State the Planar Separator Theorem.
  \item Show that \textsc{Maximum Independent Set} can be solved in
        $2^{O(\sqrt{n})}$ time on planar graphs.%
        \footnote{The pool PDF asks for $3$-coloring here. The item was changed
        to independent set; this transcription follows the amended pool.}
  \item Show that $3$-coloring is NP-hard for planar graphs.
\end{enumerate}

\subsection*{Lecture 5: Real Models of Computation}

\paragraph{Question 6. Real Models of Computation.}
This set of questions captures the learning goals of the lecture on real models
of computation.
\begin{enumerate}
  \item Explain a program on a word RAM. (Not a formal definition.)
  \item Explain a Turing machine. (Not a formal definition.)
  \item Give a comparison of advantages and disadvantages of the two models.
  \item Explain what a program on a real RAM is and give both a motivation and
        concerns about this model of computation.
  \item Show how to compute $k!$ in $O(\log^2 k)$ time on a real RAM that is
        allowed to use rounding.
  \item Show how to compute the greatest common divisor $\gcd(a,b)$ of two
        numbers in $O(\min\{\log a, \log b\})$ time on a real RAM that is
        allowed to use rounding.
  \item Show that we can find a factor of a number in polynomial time on the real
        RAM that is allowed to use rounding. You are allowed to use the answer of
        the previous two questions.
  \item Define a straight line program and \textsc{PosSLP}.
  \item Show that $\mathrm{P}^{\textsc{PosSLP}} = \mathrm{P}_{\R}$.
\end{enumerate}

\subsection*{Lecture 6: Matching}

\paragraph{Question 7.}
Consider the following problem:
\begin{probdef}{Partition into Cycles}
  \instance Undirected graph $G = (V,E)$.
  \question Can we partition $G$ into cycles, i.e., is there a collection of
  simple cycles $C_1, \ldots, C_r$ in $G$, such that each vertex belongs to
  exactly one of the cycles in the collection?
\end{probdef}
Give a polynomial time algorithm that solves this problem.

\subsection*{Lecture 7: Existential Theory of the Reals}

\paragraph{Question 8. Existential Theory of the Reals.}
This set of questions captures the learning goals of the first lecture on the
existential theory of the reals.
\begin{enumerate}
  \item Give a definition of the algorithmic problem \ETR{} and the complexity
        class $\ER$ based on this definition.
  \item Give the machine model description of $\ER$ and a sketch of a proof that
        it is equivalent with the logic based definition.
  \item Show that $\mathrm{NP} \subseteq \ER$.
  \item Define the order types and the order type problem. Give a motivation to
        study order types.
  \item Define \textsc{ETR-AM} as in the lecture and show that it is
        $\ER$-complete.
  \item Define the \textsc{Partial Order Type Problem}. Suppose that we can use
        order types to also enforce pairs of lines, defined by two points each,
        to be parallel. Show that the \textsc{Partial Order Type Problem} is
        $\ER$-complete.
\end{enumerate}

\subsection*{Lecture 8: Exact Exponential-Time Algorithms}

\paragraph{Question 9. Exact Algorithm for Exact 3-SAT.}
Suppose we have a set of Boolean variables $x_1, \ldots, x_n$, a set of clauses,
where each clause is a subset of
\[ \{x_1, \ldots, x_n, \lnot x_1, \ldots, \lnot x_n\}, \]
and each clause has at most three literals (elements). In \textsc{Exact 3-Satisfiability}, we want to
determine if there exists a truth assignment such that each clause has
\emph{exactly} one true literal.

For instance, if we have clauses $\{x_1,x_2,x_3\}$, $\{\lnot x_1, x_2, x_4\}$,
then setting $x_1$ to true, $x_2$ to false, $x_3$ to false, and $x_4$ to true is
a solution. Also, setting $x_1$ and $x_2$ to true does not lead to a solution, as
the first clause has two literals that are true.
\begin{enumerate}
  \item Show that we can create a reduction rule using polynomial time that
        transforms an instance of \textsc{Exact 3-Satisfiability} into an
        equivalent smaller instance of \textsc{Exact 3-Satisfiability} where all
        clauses have size two or three. (Smaller means, not increasing the number
        of variables or clauses.)
  \item Show that we can create a reduction rule using polynomial time that
        transforms an instance of \textsc{Exact 3-Satisfiability} into an
        equivalent smaller instance of \textsc{Exact 3-Satisfiability} where all
        clauses have size exactly three. (Smaller means, not increasing the
        number of variables or clauses.)
  \item Give an exponential-time algorithm that solves the \textsc{Exact
        3-Satisfiability} problem. Analyse the time of your algorithm. Your
        algorithm should use $\Ostar(c^n)$ time for some $c < 2$.
\end{enumerate}

\paragraph{Question 10. Bisection Width.}
Consider the following problem.
\begin{probdef}{Bisection Width}
  \instance Graph $G = (V,E)$, with $|V|$ even.
  \goal Find a partition of $V$ into two sets $V_1, V_2$, with
  $|V_1| = |V_2| = |V|/2$, such that the number of edges with one endpoint in
  $V_1$ and one endpoint in $V_2$ is as small as possible.
\end{probdef}
This problem is NP-complete, and an $\Ostar(2^n)$ algorithm is trivial. Find an
efficient exact algorithm for this problem with running time $\Ostar(c^n)$ for
some $c < 2$.

\subsection*{Lecture 9: Existential Theory of the Reals}

\paragraph{Question 11. Duality and geometric $\ER$-hardness.}
\begin{enumerate}
  \item Describe point-line duality.
  \item Show the above-below property of point-line duality.
  \item Define the cyclic order of a point with respect to a point set. Describe
        what the cyclic order corresponds to in the dual.
  \item Define geometric intersection graphs and unit disk graphs in particular.
  \item Show that deciding whether a given graph is a unit disk intersection
        graph is $\ER$-complete.
  \item Show that the \textsc{Optimal Curve Straightening} problem is
        $\ER$-complete.
\end{enumerate}

\subsection*{Lecture 10: Inclusion/Exclusion \& Approximation}

\paragraph{Question 12. Inclusion/exclusion for set cover.}
Let $U$ be any set of size $m$ ($m = |U|$), and let $F$ be a set of subsets of
$U$. A set cover $C \subseteq F$ is a set of sets that together cover all
elements in $U$, i.e., $\bigcup_{S \in C} S = U$. A minimum set cover is a set
cover using the smallest number of sets from $F$ as possible.

Give an $\Ostar(2^m)$ time and polynomial-space algorithm that counts the number
of minimum set covers of $F$.

\paragraph{Question 13. Approximation non-symmetric TSP.}
Suppose we have an instance of \textsc{Travelling Salesman Problem} that is not
symmetric but still fulfils the triangle inequality, e.g., battery usage by an
electric car in a mountainous area. We know that for every two cities $i$ and
$j$, we have that $d(i,j) \le 2 d(j,i)$, and of course, $d(j,i) \le 2 d(i,j)$.

Give a $c$-approximation algorithm for this problem that runs in polynomial
time.

\subsection*{Lecture 11: Approximation Algorithms II}

\paragraph{Question 14. Approximating Steiner Tree.}
Let $G = (V,E)$ be an undirected connected graph, let $T \subseteq V$ be a set of
vertices called the \emph{terminals}, and let $N = V \setminus T$ be the
non-terminal vertices. In the \textsc{Steiner Tree} problem, we are asked to
compute a minimum-size set of edges $S \subseteq E$ such that $S$ connects all
vertices in $T$. That is, if we let $V(S)$ be the set of endpoints of edges in
$S$, then $(V(S),S)$ is a tree with $T \subseteq V(S) \subseteq V$ connecting the
vertices in $T$, where it is allowed to use vertices from $N$ as intermediate
vertices in the tree.
\begin{enumerate}
  \item Prove that the procedure below is a $2$-approximation to \textsc{Steiner
        Tree}.

        Define the complete graph $G_T = (T,F)$ to be a graph on the terminal
        vertices $T$, with $F = \{\{u,v\} : u \in T, v \in T, u \ne v\}$ and cost
        function $c : F \to \N$, where the cost of edge $\{t_1,t_2\} \in F$
        equals the length of the shortest path from $t_1$ to $t_2$ in the
        original graph $G$. Also, define $P(\{t_1,t_2\})$ to be the set of edges
        underlying such a shortest path from $t_1$ to $t_2$ in $G$.

        Consider the following algorithm for Steiner tree: given $G = (V,E)$ and
        $T$, first compute $G_T$ and a minimum spanning tree $M \subseteq F$ on
        $G_T$ w.r.t.\ the cost function $c$, and then output the edge set
        \[ \bigcup_{m \in M} P(m). \]
        Note that a minimum spanning tree is a minimum cost edge set connecting
        \emph{all} vertices in the graph.
  \item Prove that there is no FPTAS for Steiner tree unless
        $\mathrm{P} = \mathrm{NP}$.

        You do not need part 1 to prove this, and partial points will be awarded
        if you prove the weaker statement that no additive-constant-factor
        approximation exists unless $\mathrm{P} = \mathrm{NP}$.
\end{enumerate}

\subsection*{Lecture 12: Parameterized Complexity I}

\paragraph{Question 15.}
\begin{enumerate}
  \item Define a parameterized problem, FPT, XP, and para-NP-hard.
  \item Describe the idea of a branching algorithm in general.
  \item Define \textsc{Vertex Cover} and show that \textsc{Vertex Cover} can be
        solved in FPT time using a branching algorithm.
  \item Define \textsc{Cluster Editing} and give a branching algorithm for it.
  \item Give one parameterized problem that is para-NP-complete.
  \item Define the notion of a kernel.
  \item Show that a problem is FPT if and only if it has a kernel.
  \item Show that \textsc{Vertex Cover} has a kernel of size $O(k^2)$.
  \item Define \textsc{Convex String Recoloring} and show that it has a kernel of
        size $O(k^2)$.
  \item Define the \textsc{Directed $k$-Path} problem and acyclic graphs. Show
        that \textsc{Directed $k$-Path} can be solved in polynomial time on
        acyclic graphs.
  \item Explain what a constant one-sided error algorithm is. Show that there is
        a constant one-sided error algorithm for \textsc{Directed $k$-Path} with
        running time $O(k^k \cdot n^c)$, for some suitable constant $c$.
\end{enumerate}

\subsection*{Lecture 13: Treewidth}

\paragraph{Question 16. Treewidth of grid graphs.}
The $p$ by $q$ grid graph $GR_{p \times q}$ is the graph
$(V_{p \times q}, E_{p \times q})$, with
$V_{p \times q} = \{v_{i,j} \mid 1 \le i \le p \wedge 1 \le j \le q\}$, and
$E_{p \times q} = \{\{v_{i,j}, v_{i+1,j}\} \mid 1 \le i < p \wedge 1 \le j \le q\}
\cup \{\{v_{i,j}, v_{i,j+1}\} \mid 1 \le i \le p \wedge 1 \le j < q\}$.%
\footnote{The pool PDF gives the same expression $\{v_{i,j},v_{i+1,j}\}$ twice;
the second set is clearly meant to contain the horizontal edges, and we have
transcribed it that way.}

Show, by giving a construction of a tree decomposition, that the treewidth of a
$p$ by $q$ grid graph with $p \le q$ is at most $p$.

\paragraph{Question 17. Treewidth and cliques.}
A clique in a graph $G = (V,E)$ is a subset of the vertices $C \subseteq V$ such
that for any pair of vertices $c_1, c_2 \in C$ there is an edge between the
vertices in $G$: $\{c_1,c_2\} \in E$.

Show that for any tree decomposition $T$ of $G$, $T$ must have a bag $X_i$ for
which $C \subseteq X_i$. Conclude that the treewidth of $G$ is at least
$|C| - 1$.

\subsection*{Lecture 14: Complexity for Approximation and FPT}

\paragraph{Question 18. L-reduction for Maximum Cut on Multigraphs.}
We consider the approximation complexity of the \textsc{Max Cut on Multigraphs}
problem and the \textsc{Max Not-All-Equal 3-SAT} problem.

A multigraph $G = (V,E)$ is a graph that can contain multiple identical edges.
Given a multigraph $G = (V,E)$, a cut $(V_1,V_2)$ in $G$ is a partitioning of
the vertices $V$ into two sets $V_1, V_2$ ($V_1 \cap V_2 = \emptyset$,
$V_1 \cup V_2 = V$). The value of the cut $(V_1,V_2)$ is the number of edges
$(u,v) \in E$ with one endpoint in $V_1$ and one endpoint in $V_2$. In the
\textsc{Max Cut on Multigraphs} problem, we are given a multigraph and we are
asked to compute a cut with maximum value.

In the \textsc{Max Not-All-Equal 3-SAT} problem, we are given a set of variables
$X$ and a set of clauses $C$ over the variables $X$ (clauses that can have
literals $x$ and $\lnot x$ for variables $x \in X$), where each clause consists
of \emph{exactly} two or three literals. In this problem, a clause $c \in C$ is
satisfied by a truth assignment $\varphi$, if at least one literal in $c$ is made
true by $\varphi$ and at least one literal in $c$ is made false by $\varphi$. The
question in this problem is to compute a truth assignment that satisfies the
maximum number of clauses. In this question you may use the fact that
\textsc{Max Not-All-Equal 3-SAT} is APX-complete.

Prove that \textsc{Max Cut on Multigraphs} is APX-hard by giving an L-reduction
from \textsc{Max Not-All-Equal 3-SAT}.

\subsection*{Lecture 15: Kernelization and Randomized Algorithms}

\paragraph{Question 19. Vertex Cover Kernel.}
\begin{itemize}
  \item Define \textsc{Integer Linear Programming} and show how \textsc{Vertex
        Cover} can be encoded as an Integer Linear Programming.
  \item Show that \textsc{Vertex Cover} has a linear sized kernel. (If this
        question comes up in the exam, we will ask for only part of the proof.)
  \item Define \textsc{2-List Coloring} and show that it can be solved in
        polynomial time.
  \item Using the answer to the previous question, describe a randomized constant
        one-sided error algorithm for \textsc{3-Coloring} a graph with running
        time $O(1.5^n)$.
\end{itemize}
